%% parts/p06/determinacy.tex --- Phase-4 P6 depth fill
%%
%% Register coverage: D-P7-014 and N-a-048.
%% Source accounting: v1.1 Part VII has no infinite-game or determinacy
%% development to migrate.  Every result here is new, register-driven, and
%% citation-backed; no source proof status is promoted.

\chapter{Infinite Games and Determinacy}
\label{sec:p06-determinacy}
\chaptermark{Infinite Games and Determinacy}

This chapter supplies the game-theoretic calibration separated from the
boldface pointclass definitions of the preceding chapter.  It first defines
the game data and winning condition, then distinguishes the
\(\mathsf{ZFC}\) theorems from the additional set-theoretic assertions and
their sufficient large-cardinal hypothesis.

\begin{hbdefinition}{def:p06-gale-stewart-game}{Gale--Stewart game and determinacy}
For a payoff set \(A\subseteq\omega^\omega\), the game \(G(A)\) is played in
\(\omega\) many rounds.  Player I chooses \(x(0)\), Player II chooses \(x(1)\),
and the players continue alternately, producing
\(x\in\omega^\omega\).  Player I wins exactly when \(x\in A\); otherwise
Player II wins.

A \emph{strategy} for a player is a function assigning that player's next
natural-number move to every finite position at which that player is to move.
It is \emph{winning} iff every complete play consistent with it is won by that
player.  The game, or equivalently its payoff set, is \emph{determined} iff one
of the two players has a winning strategy.  These terms name game data and a
property of that data, not a topological closure operation.  They are needed
to distinguish the ZFC theorems for open and Borel payoffs from the additional
axioms below \cite{GaleStewart1953}.
\index{Gale--Stewart game}
\index{strategy!infinite game}
\index{determined game}
\end{hbdefinition}

\begin{hbtheorem}{thm:p06-gale-stewart}{cited=GaleStewart1953}{Open and closed determinacy}
\(\mathsf{ZF}\) proves that if \(A\subseteq\omega^\omega\) is open or
closed, then \(G(A)\) is determined.  No choice principle is required.
\end{hbtheorem}

\begin{hbtheorem}{thm:p06-borel-determinacy}{cited=Martin1975}{Borel determinacy}
Every game \(G(A)\) with Borel payoff
\(A\subseteq\omega^\omega\) is determined in \(\mathsf{ZFC}\).  This is a
statement of Martin's theorem, not a local compression of its proof.
\index{Borel determinacy}
\end{hbtheorem}

\begin{hbdefinition}{def:p06-ad-pd}{Axiom of Determinacy and projective determinacy}
The \emph{Axiom of Determinacy}, \(\mathsf{AD}\), asserts that \(G(A)\) is
determined for every payoff \(A\subseteq\omega^\omega\).  \emph{Projective
determinacy}, \(\mathsf{PD}\), is the restricted assertion that \(G(A)\) is
determined whenever \(A\) belongs to a boldface projective pointclass.
These are additional set-theoretic assertions about families of games.
\(\mathsf{PD}\) is narrower than \(\mathsf{AD}\), while Borel determinacy is
already a theorem of \(\mathsf{ZFC}\); this nonredundancy is exactly what
Table~\ref{tab:p06-borel-projective} records
\cite[Section~38]{Kechris1995}.
\index{Axiom of Determinacy}
\index{projective determinacy}
\end{hbdefinition}

\begin{hbtheorem}{thm:p06-ad-versus-choice}{cited=MycielskiSteinhaus1962}{Determinacy versus Choice}
Over \(\mathsf{ZF}\), \(\mathsf{AD}\) implies that the reals cannot be
well-ordered.  Hence \(\mathsf{AD}\) contradicts the Axiom of Choice.  The
incompatibility concerns full \(\mathsf{AD}\); it does not conflict with the
\(\mathsf{ZFC}\) theorem of Borel determinacy.
\end{hbtheorem}

\begin{hbdefinition}{def:p06-woodin-calibration}{Woodin calibration}
Let \(\delta\) be a cardinal.  For \(A\subseteq V_\delta\), a cardinal
\(\kappa<\delta\) is \emph{\(\lambda\)-strong for \(A\)} when there is a
transitive class \(M\) and an elementary embedding \(j:V\to M\) such that
\[
  \operatorname{crit}(j)=\kappa,
  \qquad j(\kappa)>\lambda,
  \qquad V_\lambda\subseteq M,
  \qquad j(A)\cap V_\lambda=A\cap V_\lambda.
\]
The cardinal \(\delta\) is \emph{Woodin} iff, for every
\(A\subseteq V_\delta\), some \(\kappa<\delta\) is
\(\lambda\)-strong for \(A\) for every \(\lambda<\delta\).  This is a
conditional large-cardinal property, not an existence assertion and not the
ultrafilter definition of a measurable cardinal.  It is introduced here only
because the exact number of Woodin cardinals is the load-bearing hypothesis in
the Martin--Steel theorem; no inner-model theory is imported.
\index{Woodin cardinal}
\index{strong cardinal!relative to a set}
\end{hbdefinition}

\begin{hbtheorem}{thm:p06-martin-steel}{cited=MartinSteel1989}{Martin--Steel calibration}
Let \(n\geq1\).  If there are \(n\) Woodin cardinals below a measurable
cardinal, then games whose payoffs lie in
\(\boldsymbol{\Pi}^{1}_{n+1}\), and games with complementary
\(\boldsymbol{\Sigma}^{1}_{n+1}\) payoffs, are determined.  Consequently, if
for every positive integer \(n\) there are \(n\) Woodin cardinals below a
measurable cardinal---for example, if infinitely many Woodin cardinals lie
below one measurable cardinal---then \(\mathsf{PD}\) holds.

This is a sufficient large-cardinal hypothesis.  The theorem as stated does
not assert an equivalence, does not assert that \(\mathsf{ZFC}\) proves the
required cardinals exist, and does not identify \(\mathsf{PD}\) with
\(\mathsf{AD}\).
\end{hbtheorem}

\section{Further reading}
\label{sec:p06-further-reading}

\begin{itemize}
\item \textcite{Jech2003}: comprehensive modern coverage of axioms,
ordinals, cardinal arithmetic, stationarity, forcing, constructibility, and
large cardinals; this is the principal locator for the migrated historical
v1.1 Part VII source.
\item \textcite{Kunen2011}: an independent modern treatment of forcing and
combinatorial set theory.
\item \textcite{Halbeisen2017}: weak choice, permutation models, and a gentle
entry to forcing, including the exact atoms-to-choice-failure boundary.
\item \textcite{Bell2005}: Boolean-valued universes, generic quotients,
collapse forcing, and iterations.
\item \textcite{Kanamori2003}: historical and mathematical calibration of
the large-cardinal hierarchy.
\item \textcite{Kechris1995}: classical boldface descriptive set theory and
regularity theorems.
\item \textcite{Moschovakis2009}: descriptive set theory, games, and the
effective/classical bridge.
\item \textcite{Shelah1994Cardinal}: PCF theory and singular-cardinal
arithmetic; a pointer beyond the local treatment ceiling.
\end{itemize}

This is navigation by subject and claim boundary, not an unfiltered
bibliography dump.
